\section{Smoothness of bounded solutions}
\label{sec:bounded-hermitian-smoothness}

The continuous solutions constructed above have low regularity when the density is merely integrable. If the density is a smooth positive function of the point and the potential, every bounded solution is smooth. Differentiating that density directly would require derivatives of the solution. We instead freeze its occurrence in a smooth approximation and solve a proper exponential equation. A square of the prescribed potential in the maximum-principle function gives a trace estimate independent of its derivatives.

We first prove this estimate, retaining the matrix notation $g,h$, the traces $S,B$ and the torsion trace $\tau$ of \cref{lem:hermitian-trace-inequality}.

\begin{lemma}\label{lem:prescribed-potential-trace}
Fix $M\geq0$ and a smooth real function $G$ on $X\times[-M,M]$.
There is $C$, depending only on $\omega$, $M$ and $\|G\|_{C^2}$, with the following property. If $v,u$ are smooth, $\omega_v,\omega_u>0$, $|v|\leq M$, and
\[
 \omega_u^n=\mathrm e^{u+G(x,v)}\omega^n,
\]
then
\[
 |u|\leq\|G\|_{C^0},\qquad C^{-1}\omega\leq\omega_u\leq C\omega.
\]
\end{lemma}
\begin{proof}
The extrema of $u$ give the first bound, as in \cref{thm:smooth-proper-existence}.
Put $k=g+(v_{i\bar j})$, $S_v=\operatorname{tr}_gk$, and $D=2\|G\|_{C^0}$.
The chain rule and Cauchy--Schwarz give, for $H=u+G(x,v)$,
\[
 \Delta_gH\geq-C_0\bigl(1+S_v+|\partial v|_g^2\bigr).
\]
Indeed $\Delta_gu=S-n\geq-n$, $\Delta_gv=S_v-n$, and the mixed derivatives of $G(x,v)$ are bounded below by $-C_0(1+|\partial v|_g^2)$.
Here and below constants depend only on the declared data.

Set $r=v+M+1$, $R=2M+1$ and $s=u-\inf_Xu+1$, so
$1\leq r\leq R$ and $1\leq s\leq D+1$. For constants $A,L>0$ to be chosen, maximize
\[
 Q=\log S+Ar^2-Lu+s^{-1}
\]
on $X$. If $S<1$ at the maximum, the bounded ranges of $r,u,s$ already bound $S$ everywhere. Otherwise \cref{lem:hermitian-trace-inequality} applies there. Positivity gives
\[
 SB\geq n^2,\qquad S_v\leq S\operatorname{tr}_hk,
 \qquad |\partial v|_g^2\leq S|\partial v|_h^2.
\]
The middle inequality follows by diagonalizing $h$ relative to $g$: every eigenvalue of $h$ is at most $S$, and the diagonal entries of $k$ are positive.
Consequently,
\[
 \frac{\Delta_gH}{S}
 \geq-C_1 B-C_0\operatorname{tr}_hk-C_0|\partial v|_h^2.
\]

At the maximum,
\[
 \frac{\partial S}{S}=(L+s^{-2})\partial u-2Ar\partial v.
\]
Since $|\tau|_h^2\leq C_gB$, Young's inequality bounds the torsion term by
\[
 \left|\frac2{S^2}\operatorname{Re}(h^{i\bar j}\tau_iS_{\bar j})\right|
 \leq s^{-3}|\partial u|_h^2+A|\partial v|_h^2
   +\frac{C_g\bigl((L+1)^2(D+1)^3+AR^2\bigr)}{S^2}B.
\]
Also
\begin{align*}
 \Delta_h(Ar^2)&\geq2A\operatorname{tr}_hk-2AR B+2A|\partial v|_h^2,\\
 \Delta_h(-Lu+s^{-1})&=(L+s^{-2})(B-n)+2s^{-3}|\partial u|_h^2.
\end{align*}
Combining these inequalities with \eqref{eq:hermitian-trace-inequality} yields
\[
 \begin{split}
 0\geq\Delta_hQ\geq{}&
 \left(L-2AR-C_2-
   \frac{C_g((L+1)^2(D+1)^3+AR^2)}{S^2}\right)B-n(L+1)\\
 &+(2A-C_0)\operatorname{tr}_hk+(A-C_0)|\partial v|_h^2
    +s^{-3}|\partial u|_h^2.
 \end{split}
\]
Choose $A>C_0+1$ and then $L>2AR+C_2+2$. If the fraction in the coefficient of $B$ is at least one, $S^2$ at the maximum is bounded by its numerator. Otherwise the displayed inequality gives $B\leq n(L+1)$ there.
For $n\geq2$, the bounded determinant gives
\[
 S\leq n\mathrm e^{2\|G\|_{C^0}}B^{n-1}.
\]
For $n=1$, $S=\mathrm e^{u+G(x,v)}$ directly.
Thus $S$ is bounded at the maximum in every case. Maximality of $Q$ and the bounded ranges of $r,u,s$ give the same bound globally. The positive lower determinant bound then bounds every eigenvalue of $h$ away from zero.
\end{proof}

\Needspace{9\baselineskip}
\begin{theorem}\label{thm:bounded-hermitian-smoothness}\index{Monge--Ampere equation@Monge--Amp\`ere equation!smoothness of bounded solutions}
Let $F\in C^\infty(X\times\mathbb R,\mathbb R)$. Every bounded $u\in\PSH(X,\omega)$ satisfying
\[
 (\omega+\ddc u)^n=\mathrm e^{F(x,u)}\,\mathrm dV
\]
is smooth on $X$, and $\omega+\ddc u>0$.
\end{theorem}
\begin{proof}
The density $\mathrm e^{F(x,u)}$ is bounded and strictly positive, so \cref{cor:bounded-lp-continuity} makes $u$ continuous.
Choose smooth strictly $\omega$-plurisubharmonic $v_j\downarrow u$ by \cref{thm:bounded-exact-regularization}. They have a common bound $M$ and converge uniformly by Dini's theorem.
Define
\[
 G(x,t)=F(x,t)-t+\log\frac{\mathrm dV}{\omega^n}.
\]
By \cref{thm:smooth-proper-existence}, there is a smooth strictly $\omega$-plurisubharmonic $a_j$ with
\[
 \omega_{a_j}^n=\mathrm e^{a_j+G(x,v_j)}\omega^n.
\]
Put $\delta_j=\|G(x,v_j)-G(x,u)\|_{C^0}\to0$.
The given $u$ solves the proper equation with coefficient $G(x,u)$.
The weighted measures of $a_j-\delta_j$ and $a_j+\delta_j$ bracket its coefficient, so \cref{cor:bounded-exponential-comparison} gives
\[
 a_j-\delta_j\leq u\leq a_j+\delta_j.
\]
Hence $a_j\to u$ uniformly. By \cref{lem:prescribed-potential-trace},
$C^{-1}\omega\leq\omega_{a_j}\leq C\omega$ uniformly.
In particular $\Delta_ga_j$ is uniformly bounded. Fixed-operator elliptic estimates give uniform $W^{2,q}$ bounds for every finite $q$, so $u\in C^{1,\alpha}(X)$ for every $0<\alpha<1$.

We explain the uniformly elliptic passage to higher regularity, using the concave extension from \cref{cor:classical-higher-estimate}.
On a relatively compact coordinate chart, let $P$ send a real Hessian to its complex Hermitian part, with the coordinate normalization of $\ddc$.
All matrices $g(x)+P(D^2a_j)$ lie in a fixed compact interval $\mathcal K=\{\lambda I\leq A\leq\Lambda I\}$ of positive matrices. Define
\[
 \mathcal F(x,N)=\inf_{A\in\mathcal K}
 \left\{\log\det A+\operatorname{tr}\bigl[A^{-1}(g(x)+P(N)-A)\bigr]\right\}
       -\log\det g(x).
\]
The tangent inequality for the concave logarithmic determinant shows that this is the original operator whenever $g(x)+P(N)\in\mathcal K$.
It is concave and uniformly elliptic on real symmetric matrices: for $N'\geq0$, $P(N')\geq0$ and its trace is a fixed positive multiple of $\operatorname{tr}N'$, while $A^{\pm1}$ have uniform bounds. Its dependence on $x$ is uniformly Lipschitz on this chart.
Thus the $a_j$ are classical, and hence viscosity, solutions of
\[
 \mathcal F(x,D^2a_j)=a_j+G(x,v_j).
\]
Both sides converge uniformly except for the Hessian variable. Stability of viscosity solutions gives
$\mathcal F(x,D^2u)=u+G(x,u)$ in the viscosity sense.
The right-hand side is H\"older because $u\in C^{1,\alpha}$.
The local Evans--Krylov estimate therefore gives $u\in C^{2,\beta}$ for some $\beta>0$.
The metric bounds pass to the limit, so the original equation is now uniformly elliptic with $\omega+\ddc u>0$.
Differentiating it and applying Schauder estimates repeatedly proves $u\in C^\infty(X)$.
\end{proof}
